% !TEX root = ../3.GSM.tex
\section{Appendix}

\subsection{Milling Calculation}

\subsubsection{Determination of the Equivalent Diameter of the Material Particle}

The equivalent diameter is the diameter of a sphere having the same ratio $V/S$ as the real particle. For a spherical particle, $D_{p,\text{eq}} = D$. For a cylinder with length $L$ and diameter $D$:

The average dimensions of the rice:
\begin{itemize}
    \item Length: $L = 6$\,mm
    \item Diameter: $D = 1.5$\,mm
\end{itemize}

Considering the rice as a cylinder:
\begin{equation*}
    \frac{V}{S} = \frac{\frac{\pi D^2 L}{4}}{\pi D L + 2 \cdot \frac{\pi D^2}{4}} = \frac{D \cdot L}{4L + 2D}
\end{equation*}

\begin{equation*}
    D_{p,\text{eq}} = \frac{4V}{S} = \frac{4DL}{4L + 2D} = \frac{4 \times 1.5 \times 6}{4 \times 6 + 2 \times 1.5} = 2\,\text{mm}
\end{equation*}

Therefore: $D_{p1} = D_{p,\text{eq}} = 2$\,mm.

\subsubsection{Chart $\log\Delta\varPhi_n$ following $\log D_{pn}$}

Line equation from the theory:
\begin{equation*}
    \log\Delta\varPhi_n = (b+1)\log D_{pn} + \log K'
\end{equation*}

According to Chart 2, the diagram passes through 3 points, giving:
\begin{equation*}
    b + 1 = 2.9975 \quad \Rightarrow \quad b = 1.9975
    \qquad
    \log K' = 2.4921 \quad \Rightarrow \quad K' = 310.53
\end{equation*}

The sieve ratio:
\begin{equation*}
    r = \frac{D_{pn-1}}{D_{pn}} = \frac{1}{3}\left(\frac{0.425}{0.315} + \frac{0.315}{0.2} + \frac{0.2}{0.16}\right) = 1.39
\end{equation*}

Therefore:
\begin{equation*}
    K = \frac{K'(b+1)}{r^{b+1} - 1} = \frac{310.53 \times 2.9975}{1.39^{2.9975} - 1} = 552.93
\end{equation*}

\subsubsection{Size Distribution of Material on Sieve}

The differential equation:
\begin{equation*}
    \frac{-d\varPhi}{dD_p} = K D_p^b
    \quad \Rightarrow \quad
    \varPhi = \frac{-K}{b+1} D_p^{b+1} + C = -103.6 D_p^{2.9975} + C
\end{equation*}

When $D_p = 0.425$: $\varPhi = 0.39$ $\Rightarrow$ $C = 8.36$.

\begin{equation*}
    D_p = \sqrt[2.9975]{\frac{\varPhi - 8.36}{-103.6}}
\end{equation*}

By Bond's Law, $D_{p2}$ is the size such that 80\% of product weight passes through the sieve ($\varPhi = 0.20$):
\begin{equation*}
    D_{p2} = 0.43\,\text{mm}
\end{equation*}

\subsubsection{Milling Efficiency Calculation}

Theoretical milling capacity (dry process, $W_i = 13$\,kWh/ton):
\begin{equation*}
    T = \frac{M}{t} = \frac{200 \times 10^{-6}\,\text{ton}}{26.92/60\,\text{min}}
\end{equation*}

\begin{equation*}
    P = \frac{4}{3} \times 19 \times 13 \times \left(\frac{1}{\sqrt{0.43}} - \frac{1}{\sqrt{2}}\right) \times T = 0.13\,\text{kW}
\end{equation*}

Milling machine power:
\begin{equation*}
    P' = U \cdot I \cdot \cos\varphi = 220 \times 3.8 \times 0.8 = 668.8\,\text{W} = 0.6688\,\text{kW}
\end{equation*}

Milling efficiency:
\begin{equation}
    H = \frac{P}{P'} \times 100\% = \frac{0.13}{0.6688} \times 100\% = 19.44\%
\end{equation}

\subsection{Sifting Experiment Calculation}

\subsubsection{Determination of $F \cdot a$}

Based on Chart 1, the $\sum J_i$ curve asymptotes to:
\begin{equation*}
    F \cdot a = 25.77\,\text{g}
\end{equation*}

\subsubsection{Sifting Efficiency}
\begin{equation}
    E = \frac{J_1}{F \cdot a} \times 100\% = \frac{25}{25.77} \times 100\% = 97.01\%
\end{equation}

\subsection{Mixing Experiment Calculation}

The constituents of Soya beans (A) and Green beans (B) in an ideal mixture are:
\begin{equation*}
    C_A = \frac{a}{a+b} = \frac{3}{3+1.5} = 0.67
    \qquad
    C_B = 1 - C_A = 0.33
\end{equation*}

The standard deviation of an ideal mixture:
\begin{equation*}
    \sigma_e = \sqrt{\frac{C_A C_B}{n}}
\end{equation*}

where $N = 8$ (number of samples) and $n$ is the number of beans per sample volume.

The mixing index for each time step:
\begin{equation*}
    I_s = \sqrt{\frac{C_A C_B (N-1)}{n \cdot \sum_{i=1}^{N}(C_A - C_{iA})^2}}
\end{equation*}
